% Knowledge Graphs portfolio submission.
%
% This report follows the "Portfolio Example Structure" for the Knowledge
% Graphs course: Scenario / KG Construction / ML-based Representation /
% Logic-based Representation / Reflection. The cover pages below follow the
% pro forma and use automatic \pageref references, so that the page on which
% each learning outcome (LO) is demonstrated is always correct after
% compilation -- no manual page numbers need to be filled in.
%
% SUBMISSION NOTE: because the example structure is used, export the PDF with
% a filename ending in "-structured.pdf"
% (this file: knowledge-graphs-portfolio-structured.pdf).

\documentclass{article}
\usepackage[utf8]{inputenc}
\usepackage{amssymb}
\usepackage{booktabs}
\usepackage{geometry}
\usepackage{hyperref}
\geometry{a4paper, margin=1in}

\newcommand{\lopage}[1]{Section~\ref{#1}, page~\pageref{#1}}

\title{Hybrid Knowledge Graph Embeddings for Ranking Security-Event Technique Links:\\ A Small PyKEEN Feasibility Study}
\author{Knowledge Graphs Portfolio}
\date{\today}

\begin{document}

\maketitle

% ---------------------------------------------------------------------------
% COVER PAGES (PRO FORMA): LEARNING-OUTCOME OVERVIEW
% ---------------------------------------------------------------------------
\section*{Portfolio Cover Pages: Learning Outcomes}
The tables below follow the pro forma and map each learning outcome (LO) to
the section and page of this report where it is demonstrated. All page
references are generated automatically at compile time. For each in-scope LO,
tick the proficiency level claimed. Learning outcomes LO3 (Graph Neural
Networks) and LO10 (financial KG applications) are intentionally out of scope
for this project and are only touched on as possible extensions; no
proficiency is claimed for them.

\subsection*{Representations}
\begin{table}[h]
\centering
\begin{tabular}{p{5.2cm}p{2.6cm}p{5.4cm}}
\toprule
\textbf{Learning outcome} & \textbf{Proficiency} & \textbf{Where demonstrated} \\
\midrule
\textbf{(LO1)} Understand and apply Knowledge Graph Embeddings &
$\square$ basic \newline $\square$ exceeded &
KGE trained, compared and evaluated: \lopage{sec:ml-representation} \\
\addlinespace
\textbf{(LO2)} Understand and apply logical knowledge in KGs &
$\square$ basic \newline $\square$ exceeded &
Logic-based representation: \lopage{sec:logic-representation} \\
\addlinespace
\textbf{(LO3)} Understand and apply Graph Neural Networks &
\textit{out of scope} &
Not demonstrated; briefly noted as an extension in \lopage{sec:ml-context} \\
\addlinespace
\textbf{(LO4)} Compare different KG data models (database, semantic web, ML, data science) &
$\square$ basic \newline $\square$ exceeded &
Data model and alternatives: \lopage{sec:data-model} \\
\bottomrule
\end{tabular}
\end{table}

\subsection*{Systems}
\begin{table}[h]
\centering
\begin{tabular}{p{5.2cm}p{2.6cm}p{5.4cm}}
\toprule
\textbf{Learning outcome} & \textbf{Proficiency} & \textbf{Where demonstrated} \\
\midrule
\textbf{(LO5)} Design and implement architectures of a Knowledge Graph &
$\square$ basic \newline $\square$ exceeded &
Technologies and architecture: \lopage{sec:architecture} \\
\addlinespace
\textbf{(LO6)} Describe and apply scalable reasoning methods in KGs &
$\square$ basic \newline $\square$ exceeded &
Reasoning context for the KGE and the rules: \lopage{sec:ml-context} and \lopage{sec:logic-context} \\
\addlinespace
\textbf{(LO7)} Apply a system to create a Knowledge Graph &
$\square$ basic \newline $\square$ exceeded &
Construction pipeline: \lopage{sec:construction} \\
\addlinespace
\textbf{(LO8)} Apply a system to evolve a Knowledge Graph &
$\square$ basic \newline $\square$ exceeded &
KG evolution via KGE and via logic: \lopage{sec:ml-evolution} and \lopage{sec:logic-evolution} \\
\bottomrule
\end{tabular}
\end{table}

\subsection*{Applications}
\begin{table}[h]
\centering
\begin{tabular}{p{5.2cm}p{2.6cm}p{5.4cm}}
\toprule
\textbf{Learning outcome} & \textbf{Proficiency} & \textbf{Where demonstrated} \\
\midrule
\textbf{(LO9)} Describe and design real-world applications of KGs &
$\square$ basic \newline $\square$ exceeded &
Domain and use case: \lopage{sec:domain} \\
\addlinespace
\textbf{(LO10)} Describe financial Knowledge Graph applications &
\textit{out of scope} &
Not demonstrated; only mentioned as a transfer scenario in \lopage{sec:domain} \\
\addlinespace
\textbf{(LO11)} Apply a system to provide services through a KG &
$\square$ basic \newline $\square$ exceeded &
Service definition and outcome: \lopage{sec:service} and \lopage{sec:outcome} \\
\addlinespace
\textbf{(LO12)} Describe the connections between KGs, ML and AI &
$\square$ basic \newline $\square$ exceeded &
Connections between forms of AI: \lopage{sec:connections} \\
\bottomrule
\end{tabular}
\end{table}

\clearpage

% ---------------------------------------------------------------------------
% 1. SCENARIO
% ---------------------------------------------------------------------------
\section{Scenario}

\subsection{Domain (LO9, LO10)}
\label{sec:domain}
The domain of this knowledge graph application is \emph{enterprise security
monitoring}. A modern organisation continuously produces security-relevant
events: users authenticate to workstations and servers, processes are started
and executed by users, processes belong to recognisable software families,
files are accessed, and hosts connect to one another. A Security Operations
Centre (SOC) analyst must connect these heterogeneous facts to decide whether
a sequence of events corresponds to a known adversary behaviour, such as a
MITRE ATT\&CK technique. This is a natural knowledge-graph problem because the
value lies in the \emph{relationships} between typed entities (users, hosts,
processes, files, families, techniques), not in any single event in isolation.

A realistic real-world application built on such a graph would present an
analyst with ranked technique hypotheses for a suspicious process, together
with the provenance and the neighbouring facts that support each hypothesis.
The same architectural idea would also transfer to the financial domain
(LO10): banks likewise need to connect infrastructure, users, legal entities,
regulations and risks, and KG completion can support hidden-relationship
discovery in ownership, compliance or cyber-risk graphs. To be clear, no
financial graph is implemented in this project; the construction, embedding
and rule layers described below are simply domain-independent, and LO10
remains out of scope.

\subsection{Service (LO11)}
\label{sec:service}
The service provided by the knowledge graph is a \emph{ranked-query service}:
given a process node and the relation \texttt{suggests\_technique}, return a
ranked list of the most plausible ATT\&CK-like technique tails. This is the
analyst-facing question ``given this process, which technique is most
plausible?''. Concretely, the project asks whether a compact hybrid knowledge
graph---built from synthetic enterprise-security facts plus transparent
rule-derived hypotheses---can support link prediction for missing
\texttt{process $\rightarrow$ suggests\_technique $\rightarrow$ technique}
triples. The outcome of the project is therefore a small but complete pipeline
that constructs the graph, augments it with logic, embeds it, and answers this
structured query as a ranked recommendation. The goal is explicitly a
feasibility study on a controlled teaching graph, not an operational APT
detector, and the results in this report should be read in that light.

% ---------------------------------------------------------------------------
% 2. KG CONSTRUCTION
% ---------------------------------------------------------------------------
\section{KG Construction}
\label{sec:construction}
This part describes how the knowledge graph was constructed (LO7). The code
for this section is provided in the ZIP subfolder \texttt{2 - construction}.

\subsection{Datasets}
\label{sec:datasets}
The dataset is a deterministic, self-constructed synthetic enterprise-security
knowledge graph with 1{,}232 labelled triples, 245 entities and 9 relations.
The observed-style component contains 944 triples representing pseudonymous
users, workstations, servers, processes, file kinds and asset roles. A
transparent rule layer derives a further 288 facts. The graph labels are
semantically grounded in MITRE ATT\&CK technique identifiers, but the events
and labels are generated locally rather than collected from real incidents:
this is a synthetic teaching dataset, not real telemetry.

The dataset is self-constructed and therefore included directly in the ZIP
file (folder \texttt{2 - construction}, with \texttt{base\_triples.tsv},
\texttt{derived\_rule\_triples.tsv} and \texttt{all\_triples.tsv}); it is also
fully reproducible by running the generator. The original project plan was to
use the full LANL ``Comprehensive, Multi-Source Cyber-Security Events''
dataset (\url{https://lanl.ma.ic.ac.uk/data/cyber1/}), which contains
authentication, process, network-flow, DNS and red-team-labelled events. That
dataset exceeded 60~GB, and processing at that scale on available hardware
(and on Google Colab) led to intractable runtimes and incomplete outputs. An
experimental repository documents this effort
(\url{https://github.com/arslansmajevic/knowledge-graphs-project}). As a
recovery path, this project uses a small, reproducible, interpretable
synthetic dataset that demonstrates the same hybrid architecture and
learning-outcome coverage at a tractable scale.

\subsection{Technologies and Architecture (LO5)}
\label{sec:architecture}
The system uses Python, with PyKEEN for knowledge-graph embedding and
evaluation and a small dependency-free Python module for the symbolic rule
layer; data are stored as plain labelled triples in TSV files. The
architecture is deliberately modular, mirroring how a production KG system
separates raw telemetry, symbolic rules and ML services:
\begin{quote}
synthetic event facts (observed-style) $\rightarrow$ transparent rules
(derived facts) $\rightarrow$ hybrid KG (triples) $\rightarrow$ PyKEEN KGE
(ranking / evaluation).
\end{quote}
Each stage is an independent component: \texttt{generate\_data.py} creates the
observed-style triples, \texttt{rules.py} derives the symbolic facts with
provenance, and \texttt{train.py} splits the KG, trains the embedding models
and produces ranked output. Keeping ingestion, rule derivation, triple
storage, embedding training and ranked output as distinct stages is what makes
the pipeline an architecture rather than a single script. In a real
deployment, raw telemetry would live outside the embedding model, and only
selected, versioned KG facts would be embedded.

\subsection{Construction}
\label{sec:construction-detail}
The graph contains users, workstations, servers, processes, file kinds,
process families, security roles and ATT\&CK-like technique nodes. Three
concrete construction examples illustrate how typed nodes and edges are built
from the datasets:
\begin{enumerate}
    \item \textbf{Observed edge.} An observed-style log fact yields an
    \texttt{authenticates\_to} edge between a user node and a host node, e.g.
    \texttt{(user\_06, authenticates\_to, workstation\_12)}, materialising the
    raw telemetry as a typed labelled triple.
    \item \textbf{Typing edge.} A process is connected to its software family
    with \texttt{has\_family}, e.g.
    \texttt{(process\_000, has\_family, family\_powershell)}. This edge later
    becomes the antecedent that the rule layer reasons over.
    \item \textbf{Derived edge.} The rule layer creates a
    \texttt{suggests\_technique} edge between a process and an ATT\&CK-like
    technique node, e.g.
    \texttt{(process\_000, suggests\_technique,\\
    attack\_T1059\_command\_and\_scripting)},
    together with a \texttt{has\_provenance} fact recording that it was
    produced by the rule rather than observed directly.
\end{enumerate}
Observed relations include \texttt{authenticates\_to}, \texttt{runs},
\texttt{executed\_by}, \texttt{has\_family}, \texttt{accesses},
\texttt{connects\_to} and \texttt{has\_role}; the derived relations are
\texttt{suggests\_technique} and \texttt{has\_provenance}.

% ---------------------------------------------------------------------------
% 3. ML-BASED REPRESENTATION
% ---------------------------------------------------------------------------
\section{ML-based Representation}
\label{sec:ml-representation}
This part demonstrates Knowledge Graph Embeddings (LO1). The code for this
section is provided in the ZIP subfolder \texttt{3 - ML}.

\subsection{The ML Representation and Training (LO1)}
\label{sec:ml-training}
The ML-based representation is a knowledge-graph embedding that maps every
entity and relation to a continuous vector. Triples were randomly split
80/10/10 into training, validation and test sets using split seed 42. The two
models were TransE and ComplEx, trained with PyKEEN using Adam optimisation,
learning rate 0.01 and batch size 256. The main metrics are filtered mean
reciprocal rank (MRR), Hits@1, Hits@3, Hits@10 and mean rank; the project also
reports an application-aligned \texttt{technique\_tail\_type\_constrained}
metric that ranks only candidate technique tails for
\texttt{suggests\_technique} (and is therefore easier than global link
prediction; the two must not be compared numerically).

\paragraph{Five example representations.} TransE represents the relation
\texttt{suggests\_technique} as a translation vector, so processes from the
same family are moved in a consistent direction toward the same technique
node, for example: (1) \texttt{process\_000} $\to$
\texttt{attack\_T1059\_command\_and\_scripting}; (2) \texttt{process\_009}
embedded near \texttt{process\_000} because they share
\texttt{family\_powershell}; (3) the relation \texttt{has\_family} as a
separate translation linking processes to family nodes; (4) the entity
\texttt{workstation\_12} embedded close to the users that
\texttt{authenticate\_to} it; (5) the technique node
\texttt{attack\_T1003\_os\_credential\_dumping} positioned as the shared tail
of all credential-dump-family processes.

\paragraph{True and false positives.} A true-positive ranking is
\texttt{(process\_000, suggests\_technique,
attack\_T1059\_command\_and\_scripting)}, which the model ranks at the top of
the candidate list and which the rule layer also supports. A false-positive
ranking is a non-technique entity such as \texttt{process\_009} or
\texttt{file\_spreadsheet\_09} appearing high in the unconstrained candidate
list for the same query---these are structurally plausible neighbours but are
not valid technique tails, which is exactly why the type-constrained metric is
reported separately.

\paragraph{Results.} The experiments cover a fair model comparison (Run~A),
and sensitivity to training duration (Run~B), embedding dimension (Run~C), a
final held-out test (Run~D) and seed robustness (Run~E). All numbers below are
taken directly from the \texttt{metrics\_comparison.csv} files produced by the
pipeline.

\begin{table}[h]
\centering
\begin{tabular}{lrrrrr}
\toprule
\textbf{Model} & \textbf{MRR} & \textbf{Hits@1} & \textbf{Hits@3} & \textbf{Hits@10} & \textbf{Mean Rank} \\
\midrule
TransE & 0.244 & 0.110 & 0.301 & 0.524 & 32.00 \\
ComplEx & 0.027 & 0.004 & 0.008 & 0.041 & 100.81 \\
\bottomrule
\end{tabular}
\caption{Run~A. Validation results for TransE and ComplEx, both trained for 80 epochs with dimension 32.}
\end{table}

TransE clearly outperformed ComplEx under these conditions. This is a finding
about this graph and this training budget, not a universal claim. The
synthetic rules create simple translational patterns (same-family processes
point to the same technique), which TransE models efficiently; ComplEx is more
expressive but that capacity did not help within the same 80-epoch budget.

\begin{table}[h]
\centering
\begin{tabular}{rrrrrr}
\toprule
\textbf{Epochs} & \textbf{MRR} & \textbf{Hits@1} & \textbf{Hits@3} & \textbf{Hits@10} & \textbf{Mean Rank} \\
\midrule
20 & 0.159 & 0.041 & 0.195 & 0.390 & 52.14 \\
80 & 0.244 & 0.110 & 0.301 & 0.524 & 32.00 \\
200 & 0.361 & 0.232 & 0.419 & 0.602 & 20.76 \\
\bottomrule
\end{tabular}
\caption{Run~B. Validation results for TransE (dimension 32) at different epoch counts.}
\end{table}

Validation improved from 20 to 80 to 200 epochs, indicating the 80-epoch model
had not fully converged. If the final model is selected strictly on
validation, the 200-epoch configuration is the most consistent next choice.

\begin{table}[h]
\centering
\begin{tabular}{rrrrrr}
\toprule
\textbf{Dimension} & \textbf{MRR} & \textbf{Hits@1} & \textbf{Hits@3} & \textbf{Hits@10} & \textbf{Mean Rank} \\
\midrule
16 & 0.299 & 0.179 & 0.333 & 0.549 & 26.16 \\
32 & 0.244 & 0.110 & 0.301 & 0.524 & 32.00 \\
64 & 0.213 & 0.057 & 0.305 & 0.528 & 32.31 \\
\bottomrule
\end{tabular}
\caption{Run~C. Validation results for TransE (80 epochs) at different embedding dimensions.}
\end{table}

The 16-dimensional model had the best global validation MRR, even though the
64-dimensional model had the highest type-constrained technique-tail MRR. The
graph is small and regular, so a compact embedding suffices; more capacity is
not automatically better, and global and task-specific metrics must be
interpreted separately.

\begin{table}[h]
\centering
\begin{tabular}{lrrrrr}
\toprule
\textbf{Model} & \textbf{MRR} & \textbf{Hits@1} & \textbf{Hits@3} & \textbf{Hits@10} & \textbf{Mean Rank} \\
\midrule
TransE & 0.240 & 0.101 & 0.286 & 0.532 & 30.74 \\
\bottomrule
\end{tabular}
\caption{Run~D. Final test result for the available TransE 80-epoch, 32-dimensional configuration.}
\end{table}

The test MRR of 0.240 is close to the validation MRR of 0.244 for the same
configuration, so the validation estimate was not wildly optimistic. Hits@10
of 0.532 means the true held-out triple was ranked in the top 10 in just over
half of the evaluated cases---meaningful for a feasibility study, but
insufficient for an operational system.

\begin{table}[h]
\centering
\begin{tabular}{rrrrrr}
\toprule
\textbf{Seed} & \textbf{MRR} & \textbf{Hits@1} & \textbf{Hits@3} & \textbf{Hits@10} & \textbf{Mean Rank} \\
\midrule
7 & 0.256 & 0.097 & 0.347 & 0.577 & 25.10 \\
42 & 0.240 & 0.101 & 0.286 & 0.532 & 30.74 \\
1337 & 0.296 & 0.177 & 0.327 & 0.556 & 29.27 \\
\bottomrule
\end{tabular}
\caption{Run~E. Test robustness across three model seeds for TransE (80 epochs, dimension 32).}
\end{table}

The three test MRR values (0.256, 0.240, 0.296) have a mean of about 0.264 and
a standard deviation of about 0.029, indicating moderate stability across
seeds.

\subsection{Evolving the KG with the ML Representation (LO8)}
\label{sec:ml-evolution}
The ML representation is used to \emph{complete} the knowledge graph. Link
prediction proposes missing \texttt{suggests\_technique} edges: a held-out
true triple is the KG-completion task, and a top-ranked candidate is a new
hypothesis edge that could be added to the graph. The effect on evolution is
therefore \emph{adding} a candidate edge. Because KGE scores are rankings
rather than calibrated probabilities, each proposed edge is treated as a
candidate update for analyst review, not as confirmed truth.

\subsection{Context and Limitations (LO6)}
\label{sec:ml-context}
The KGE performs \emph{approximate} reasoning: it ranks plausible missing
links rather than deriving them deterministically, which is the scalable
completion counterpart to the exact rule reasoning of
Section~\ref{sec:logic-context}. On scaling: \emph{volume}---PyKEEN's filtered
evaluation and negative sampling grow with the number of entities and triples,
so a 60~GB log graph would need batching, sampling and likely GPU training.
\emph{Variety}---adding new relation types is cheap for the triple model but
dilutes the regular translational pattern that made TransE effective here.
\emph{Veracity}---with little or noisy training data, embedding quality
degrades quickly, as the seed-robustness spread already hints. In our concrete
scenario, the model only ever sees four technique families, so its apparent
strength would not transfer to the full ATT\&CK matrix without far more, and
more diverse, training evidence. Two further honest caveats: the random triple
split is appropriate for a teaching benchmark but could leak temporal
information in a real log setting (a real extension needs a chronological
split), and the ranking metrics test recovery of withheld triples, not alert
precision or recall on independently labelled incidents. Graph Neural
Networks (LO3) were not applied in this project; a relational GNN over the
same graph would be a natural extension, but no claim is made about it here.

% ---------------------------------------------------------------------------
% 4. LOGIC-BASED REPRESENTATION
% ---------------------------------------------------------------------------
\section{Logic-based Representation}
\label{sec:logic-representation}
This part demonstrates logical knowledge in KGs (LO2). The code for this
section is provided in the ZIP subfolder \texttt{4 - logic}.

\subsection{The Logic Representation (LO2)}
\label{sec:logic-detail}
The logic-based representation is a small, inspectable, deterministic
Datalog-style rule set, kept deliberately separate from model training. The
core rule is:
\begin{quote}
\texttt{has\_family(Process, Family)} and
\texttt{family\_maps\_to\_technique(Family, Technique)} imply
\texttt{suggests\_technique(Process, Technique)}.
\end{quote}
Five concrete examples of the logic-based representation:
\begin{enumerate}
    \item \texttt{family\_powershell} $\mapsto$
    \texttt{attack\_T1059\_command\_and\_scripting}.
    \item \texttt{family\_remote\_desktop} $\mapsto$
    \texttt{attack\_T1021\_remote\_services}.
    \item \texttt{family\_downloader} $\mapsto$
    \texttt{attack\_T1105\_ingress\_tool\_transfer}.
    \item \texttt{family\_credential\_dump} $\mapsto$
    \texttt{attack\_T1003\_os\_credential\_dumping}.
    \item For every derived \texttt{suggests\_technique} fact the rule also
    \emph{creates} a new \texttt{has\_provenance} edge
    (\texttt{rule\_family\_to\_attack\_technique}), so the symbolic reason for
    a suggestion is stored explicitly in the graph.
\end{enumerate}
Example~5 illustrates a rule that \emph{creates} new nodes/edges in the graph.
To be explicit about scope: the implemented rule set is a single
non-recursive step and does not implement full recursive or existential
Datalog reasoning. A natural recursive extension is described in
Section~\ref{sec:logic-context} (following \texttt{connects\_to} chains for
lateral-movement detection across arbitrarily many edges).

\subsection{Evolving the KG with the Logic Representation (LO8)}
\label{sec:logic-evolution}
The logic-based representation is used to \emph{complete} the knowledge graph
by \emph{adding} new facts: each firing of the rule adds a
\texttt{suggests\_technique} edge and a \texttt{has\_provenance} edge that
were not present in the observed-style data. Unlike the KGE, this addition is
deterministic and fully auditable---given the same observed facts, the same
derived facts are produced---which is why the derived triples are stored with
provenance and used as transparent ground for the embedding layer.

\subsection{Context and Limitations (LO6)}
\label{sec:logic-context}
The rule engine performs \emph{deterministic, exact} reasoning, which is the
complement to the KGE's approximate completion. On scaling within our
scenario: \emph{volume}---deterministic forward chaining over a single
non-recursive rule is linear in the number of \texttt{has\_family} facts and
scales well, but a recursive lateral-movement rule over \texttt{connects\_to}
chains could grow combinatorially and would need stratification or bounded
recursion. \emph{Variety}---the current rule covers only four families, so
extending it to the full ATT\&CK matrix means authoring and maintaining many
more mappings. \emph{Veracity}---the rule trusts its \texttt{has\_family}
input, so noisy or spoofed family labels would propagate directly into wrong
technique suggestions. A concrete improvement is to add a recursive rule that
follows host-to-host \texttt{connects\_to} edges, plus richer provenance and
temporal conditions, so the logic layer can express multi-hop attack chains
rather than a single family-to-technique step.

% ---------------------------------------------------------------------------
% 5. REFLECTION
% ---------------------------------------------------------------------------
\section{Reflection}
\label{sec:reflection}
The supporting material for this part is provided in the ZIP subfolder
\texttt{5 - reflection}.

\subsection{Outcome of the Service (LO11)}
\label{sec:outcome}
The service formulated in Section~\ref{sec:service} was delivered: the
pipeline answers ``given this process, which technique is most plausible?'' as
a ranked list, and the final TransE test run places the true held-out
technique in the top 10 for just over half of the evaluated cases (Hits@10
$= 0.532$). As a ranked-query service this is a working, if modest, outcome:
it demonstrates feasibility on a controlled graph. It is not yet a deployable
analyst tool because the scores are rankings, not calibrated risk values, and
the technique labels are rule-derived rather than independently observed.
Human review therefore remains necessary. The most consistent improvement to
the service is to run the final test for the strongest validation
configuration (TransE, 200 epochs, dimension 32) and to surface provenance and
neighbouring facts alongside each ranked hypothesis.

\subsection{Data Model and Alternatives (LO4)}
\label{sec:data-model}
The chosen data model is a labelled-triple representation. The generator
naturally resembles a \emph{property graph} from the database community,
because security objects are typed entities connected by labelled edges;
PyKEEN then consumes the same content as labelled triples, which is close to
the \emph{RDF-style} statement view from the semantic-web community. Compared
with a pure relational model, the triple/property-graph view makes multi-hop
relationship queries (the core of the scenario) far more natural, at the cost
of weaker native support for tabular aggregation. Compared with a
feature-matrix view from the machine-learning and data-science communities,
the graph view preserves explicit relations and provenance that a flattened
feature table would discard. The triple model was chosen because the project's
value lies in relationships, and because it is the shared substrate that both
the symbolic rules and the embedding model can operate on without conversion.
This is a modelling choice suited to this project, not a claim that one data
model is universally better.

\subsection{Connections between Forms of AI (LO12)}
\label{sec:connections}
The project combines three forms of knowledge: symbolic rules, graph data
modelling and neural embeddings. The interaction in this scenario is one-way
and sequential: the logic layer produces transparent
\texttt{suggests\_technique} and \texttt{has\_provenance} facts, and the
embedding layer then generalises over the whole graph, including those facts,
to rank missing links. This interaction could be improved by making it
bidirectional---for example, using rule provenance to mask rule-derived edges
during evaluation so the KGE is measured on genuinely new generalisation, or
using high-confidence KGE predictions to propose candidate new rules. It would
be partly possible to learn the logical mappings automatically (e.g. rule
mining over observed family/technique co-occurrences), and the KGE's accuracy
could in turn be improved by logical constraints that forbid type-invalid
tails (the false positives noted in Section~\ref{sec:ml-training}). Because
the targets are currently rule-derived, the embedding partly relearns the
rule; the clearest path to a genuinely hybrid system is to introduce
independently sourced labels so that the logic and ML layers contribute
complementary, rather than overlapping, knowledge. Neither component alone
makes the system a complete AI-based security product.

\end{document}
